\medskip\noindent\textbf{Part 1 (15 points): The Time Problem}

\begin{description}
\item[(a)] Local sidereal time:
  \[ \mathrm{LST} = \alpha + H \]
  \hfill(1 point)

  where $\alpha$ is the star's Right Ascension.

  The time difference:
  \[ \Delta T = \mathrm{LST} - T_R \]
  \hfill(1 point)

  Star A:
  \begin{align*}
    \mathrm{LST} &= 17^{\mathrm{h}}35^{\mathrm{m}}35.55^{\mathrm{s}}
      + 2^{\mathrm{h}}02^{\mathrm{m}}08.95^{\mathrm{s}}
      = 19^{\mathrm{h}}37^{\mathrm{m}}44.50^{\mathrm{s}} \\
    \Delta T_{\mathrm{A}} &= 19^{\mathrm{h}}37^{\mathrm{m}}44.50^{\mathrm{s}}
      - 19^{\mathrm{h}}36^{\mathrm{m}}53.69^{\mathrm{s}}
      = \mathbf{50.81\,s}
  \end{align*}
  \hfill(0.5 points)

  Star B:
  \begin{align*}
    \mathrm{LST} &= 19^{\mathrm{h}}59^{\mathrm{m}}52.48^{\mathrm{s}}
      + 1^{\mathrm{h}}47^{\mathrm{m}}50.38^{\mathrm{s}}
      = 21^{\mathrm{h}}47^{\mathrm{m}}42.86^{\mathrm{s}} \\
    \Delta T_{\mathrm{B}} &= 21^{\mathrm{h}}47^{\mathrm{m}}42.86^{\mathrm{s}}
      - 21^{\mathrm{h}}46^{\mathrm{m}}52.11^{\mathrm{s}}
      = \mathbf{50.75\,s}
  \end{align*}
  \hfill(0.5 points)

  Star C:
  \begin{align*}
    \mathrm{LST} &= 01^{\mathrm{h}}39^{\mathrm{m}}52.31^{\mathrm{s}}
      - 02^{\mathrm{h}}00^{\mathrm{m}}30.38^{\mathrm{s}} + 24^{\mathrm{h}}
      = 23^{\mathrm{h}}39^{\mathrm{m}}21.93^{\mathrm{s}} \\
    \Delta T_{\mathrm{C}} &= 23^{\mathrm{h}}39^{\mathrm{m}}21.93^{\mathrm{s}}
      - 23^{\mathrm{h}}38^{\mathrm{m}}31.61^{\mathrm{s}}
      = \mathbf{50.32\,s}
  \end{align*}
  \hfill(0.5 points)

  Star D:
  \begin{align*}
    \mathrm{LST} &= 01^{\mathrm{h}}41^{\mathrm{m}}38.84^{\mathrm{s}}
      - 01^{\mathrm{h}}37^{\mathrm{m}}13.18^{\mathrm{s}}
      = 00^{\mathrm{h}}04^{\mathrm{m}}25.66^{\mathrm{s}} \\
    \Delta T_{\mathrm{D}} &= 00^{\mathrm{h}}04^{\mathrm{m}}25.66^{\mathrm{s}}
      - 00^{\mathrm{h}}03^{\mathrm{m}}35.22^{\mathrm{s}}
      = \mathbf{50.44\,s}
  \end{align*}
  \hfill(0.5 points)

  \[
  \rightarrow \Delta T_{\mathrm{mean}}
  = \frac{50.81 + 50.75 + 50.32 + 50.44}{4} = \mathbf{50.58\,s}
  \]
  \hfill(1 point)
  \[
  \rightarrow \sigma_T
  = \left(\frac{0.23^2 + 0.17^2 + 0.26^2 + 0.14^2}{4}\right)^{0.5}
  = \mathbf{0.21\,s}
  \]
  \hfill(1 point)

\item[(b)] Stars A and B were observed in the West (setting).
  \[
  \Delta T_{\mathrm{W}}
  = \frac{\Delta T_{\mathrm{A}} + \Delta T_{\mathrm{B}}}{2}
  = \frac{50.81 + 50.75}{2} = \mathbf{50.78\,s}
  \]
  \hfill(1 point)

  Stars C and D were observed in the East (rising).
  \[
  \Delta T_{\mathrm{E}}
  = \frac{\Delta T_{\mathrm{C}} + \Delta T_{\mathrm{D}}}{2}
  = \frac{50.32 + 50.44}{2} = \mathbf{50.38\,s}
  \]
  \hfill(1 point)

  The discrepancy is:
  \[
  \Delta T_{\mathrm{W}} - \Delta T_{\mathrm{E}} = 50.78 - 50.38
  = \mathbf{0.40\,s}
  \]
  \hfill(1 point)

\item[(c)] The time difference:
  \[ \Delta T = \mathrm{LST} - T_R \]
  The error on time measurement:
  \[
  \mathrm{d}(\Delta T) = \mathrm{d}(\mathrm{LST}) = \mathrm{d}(\alpha + H)
  = \mathrm{d}H
  \]
  \hfill(1.5 points)

  Use the Sine-Cosine Rule for the triangle Pole-Zenith-Star (PZS), we have:
  \[
  \cos z = \sin\phi \sin\delta + \cos\phi \cos\delta \cos H \quad (1)
  \]
  \hfill(1.5 points)

  where $\phi$ is the latitude of the observatory and $\delta$ is the
  declination of the considered star.

  Differentiate equation (1):
  \[
  -\sin z\,\mathrm{d}z = -\cos\phi \cos\delta \sin H\,\mathrm{d}H
  \]
  Hence:
  \[
  \mathrm{d}(\Delta T) = \mathrm{d}H
  = \frac{\sin z\,\mathrm{d}z}{\cos\phi \cos\delta \sin H}
  \]
  \hfill(1 point)

  For East stars:
  \[ H < 0 \rightarrow \sin H < 0 \]
  \hfill(0.5 points)

  For West stars:
  \[ H > 0 \rightarrow \sin H > 0 \]
  \hfill(0.5 points)

  Thus, the same zenith-distance error $(\mathrm{d}z)$ produces hour-angle
  corrections $(\mathrm{d}H)$ of the same magnitude but opposite sign.
  Averaging East and West observations cancels much of the systematic error
  caused by the astrolable:
  \[
  \frac{1}{2}\left[\mathrm{d}(\Delta T_{\mathrm{W}})
    + \mathrm{d}(\Delta T_{\mathrm{E}})\right]
  < \left|\mathrm{d}(\Delta T_{\mathrm{W,E}})\right|
  \]
  \hfill(1 point)
\end{description}

\medskip\noindent\textbf{Part 2 (3 points): Atmospheric Refraction}

Atmospheric refraction in radians:
\[
R = \frac{7.888 \times 10^{-5} \times 1011.0}
         {[(273.0 + 27.0) \times \tan 60^\circ]}
  = 1.5347 \times 10^{-4}\,\mathrm{rad}
\]
\hfill(1 point)

Convert to arcseconds:
\[
R = 1.5347 \times 10^{-4} \times 206264.8 = \mathbf{31.66''}
\]
\hfill(1 point)

The true zenith distance:
\[
z_{\mathrm{true}} = 30^\circ 00' 00'' + 31.66'' = \mathbf{30^\circ 00' 31.66''}
\]
\hfill(1 point)

\emph{(The `+' sign because the atmospheric refraction acts like a lens,
apparent altitude: $h_{\mathrm{app}} = h_{\mathrm{true}} + R \rightarrow
(90 - z_{\mathrm{app}}) = (90 - z_{\mathrm{true}}) + R \rightarrow
z_{\mathrm{true}} = z_{\mathrm{app}} + R$)}

\medskip\noindent\textbf{Part 3 (15 points): Yersin's Lines of Position}

\begin{description}
\item[(a)] Use the Law of Cosine for the triangle PZS:
  \[
  \sin\delta = \sin\phi \cos z + \cos\phi \sin z \cos A
  \]
  \hfill(1.5 points)
  \[
  \rightarrow \cos\phi \sin z \cos A = \sin\delta - \sin\phi \cos z
  \]
  Substitute $\cos z$ into the right side of the above equation:
  \begin{align*}
    \cos\phi \sin z \cos A
      &= \sin\delta - \sin\phi\,(\sin\phi \sin\delta
         + \cos\phi \cos\delta \cos H) \\
    \rightarrow \cos\phi \sin z \cos A
      &= \sin\delta \cos^2\phi - \sin\phi \cos\phi \cos\delta \cos H \\
    \rightarrow \sin z \cos A
      &= \cos\phi \sin\delta - \sin\phi \cos\delta \cos H \quad (2)
  \end{align*}
  \hfill(1 point)

  From the Spherical Law of Sines:
  \[
  \frac{\sin H}{\sin z} = -\frac{\sin A}{\cos\delta}
  \]
  \hfill(1.5 points)

  Rearranging:
  \[
  \sin z \sin A = -\cos\delta \sin H \quad (3)
  \]
  \emph{(the minus sign because $H$ is Westward and $A$ is Eastward)}

  Substitute the identities (2) and (3) into the derivative of the zenith
  distance equation (in Hints):
  \begin{align*}
    -\sin z\,\mathrm{d}z
      &= (\sin z \cos A)\,\mathrm{d}\phi
         + (\sin z \sin A)\cos\phi\,\mathrm{d}H \\
    \mathrm{d}z &= -\sin A\,(\cos\phi\,\mathrm{d}H) - \cos A\,\mathrm{d}\phi
      \quad \textbf{(Q.E.D)}
  \end{align*}
  \hfill(1 point)

\item[(b)]
  % FIGURE NEEDED: 2026_DASOL_D02 crop page 6, hand-drawn Graph D02.3b: the four LOPs A, B, C, D plotted on the (x, y) plane in arcsec, with x from -10 to 30 and y from 0 to 50, intersecting near (8.5, 27)

  \textbf{Plotting}: Draw the four LOPs $x\sin(A) + y\cos(A) = p$ using the
  supplied azimuths and intercepts which forms an intersection region around
  the center $(8.5'', 27'')$.

  Line A: $(0, -2)$ $(10, 36)$ \hfill(1.5 points)

  Line B: $(0, 44)$ $(21, 0)$ \hfill(1.5 points)

  Line C: $(0, 53)$ $(19, 0)$ \hfill(1.5 points)

  Line D: $(0, 14)$ $(-10, 0)$ \hfill(1.5 points)

  $x_0$ and $y_0$ values are within the acceptable ranges:
  \[
  7'' \le x_0 \le 10'', \qquad 26'' \le y_0 \le 30''
  \]
  \hfill(4 points)
\end{description}

\medskip\noindent\textbf{Part 4 (17 points): Recovering the Location of Yersin Observatory and Error Analytics}

\begin{description}
\item[(a)] Latitude correction:
  \begin{align*}
    y_0 &= \mathrm{d}\phi = \phi_{\mathrm{true}} - \phi_0 \\
    \rightarrow \phi_{\mathrm{true}} &= \phi_0 + y_0
  \end{align*}
  Therefore:
  \[
  \phi_{\mathrm{true}} = 12^\circ 15' 00'' + 27.20''
  = \mathbf{12^\circ 15' 27.20''}
  \]
  \hfill(1 point)

  Using $x_0 = 8.50''$ and $\phi_0 = 12^\circ 15'$ ($\cos\phi_0 \approx 0.977$)
  gives:
  \[
  \mathrm{d}T_{\mathrm{corr}} = \frac{x_0}{15\cos\phi_0}
  = \frac{8.50}{15 \times 0.977} = \mathbf{0.58\,s}
  \]
  \hfill(2 points)

  Hence:
  \[
  \Delta T_{\mathrm{true}} = \Delta T_0 + \mathrm{d}T_{\mathrm{corr}}
  = 50.00 + 0.58 = \mathbf{50.58\,s}
  \]
  \hfill(1 point)

\item[(b)] From the result of question c) in Part 1, we have:
  \[
  \mathrm{d}(\Delta T) = \frac{\sin z\,\mathrm{d}z}{\cos\phi \cos\delta \sin H}
  \]
  \hfill(1 point)

  Using the Spherical Law of Sines:
  \[
  \frac{\sin H}{\sin z} = -\frac{\sin A}{\cos\delta}
  \]
  \hfill(1 point)
  \[
  \rightarrow \mathrm{d}(\Delta T) = -\frac{\mathrm{d}z}{\cos\phi \sin A}
  \]
  \hfill(1 point)

  For the Eastern star (the average azimuth $\sim 90^\circ$):
  \[
  \mathrm{d}(\Delta T_{\mathrm{E}}) = -\frac{\mathrm{d}z}{\cos\phi}
  \]
  \hfill(1 point)

  For the Western stars (the average azimuth $\sim 270^\circ$):
  \[
  \mathrm{d}(\Delta T_{\mathrm{W}}) = \frac{\mathrm{d}z}{\cos\phi}
  \]
  \hfill(1 point)

  Combining these results yields:
  \[
  \left|\mathrm{d}(\Delta T_{\mathrm{W}}) - \mathrm{d}(\Delta T_{\mathrm{E}})\right|
  = \frac{2\,\mathrm{d}z}{\cos\phi}
  \]
  \hfill(1 point)

  Convert into seconds:
  \[
  \Delta T_{\mathrm{W}} - \Delta T_{\mathrm{E}}
  = \frac{2\,\mathrm{d}z}{15\cos\phi}
  \]
  \hfill(1 point)

  Therefore:
  \[
  \mathrm{d}z = \frac{15}{2}\,(\Delta T_{\mathrm{W}} - \Delta T_{\mathrm{E}})
  \cos\phi_0 = \mathbf{2.93''}
  \]
  \hfill(1 point)

\item[(c)]
  \[
  \varepsilon_{\mathrm{A}} = \left|8.50''\sin(285.0^\circ)
    + 27.20''\cos(285.0^\circ) + 0.40''\right| = \mathbf{0.77''}
  \]
  \hfill(0.5 points)
  \[
  \varepsilon_{\mathrm{B}} = \left|8.50''\sin(245.0^\circ)
    + 27.20''\cos(245.0^\circ) + 18.77''\right| = \mathbf{0.43''}
  \]
  \hfill(0.5 points)
  \[
  \varepsilon_{\mathrm{C}} = \left|8.50''\sin(70.0^\circ)
    + 27.20''\cos(70.0^\circ) - 18.08''\right| = \mathbf{0.79''}
  \]
  \hfill(0.5 points)
  \[
  \varepsilon_{\mathrm{D}} = \left|8.50''\sin(125.0^\circ)
    + 27.20''\cos(125.0^\circ) + 8.15''\right| = \mathbf{0.49''}
  \]
  \hfill(0.5 points)

\item[(d)]
  \[
  d = \sqrt{x_0^2 + y_0^2} = \sqrt{8.50''^2 + 27.20''^2} = \mathbf{28.50''}
  \]
  \hfill(2 points)

  $1''$ is roughly $31\,\mathrm{m}$, therefore,
  $d = 28.50'' \times 31\,\mathrm{m}/1'' \approx \mathbf{883\,m}$.
  \hfill(1 point)
\end{description}
